\chapter*{Résumé}
\addcontentsline{toc}{chapter}{Résumé}

\paragraph{Objet.} La phase 3.1 consolide la version 0.7 du modèle sur six points : le paramètre qui transmet la
demande au PIB potentiel (dont dépendait la dette de long terme de S3, de $-9$ à $+43$ points), la validation sur le
passé (backtest 2010-2025), l'incidence des prélèvements sur les entreprises, la distribution des effets entre ménages
après propagation macroéconomique, la cohérence entre modules et la réévaluation des stratégies S1 à S5 (définitions
inchangées). Toutes les simulations citées ont été exécutées ; les résultats ne sont ni des prévisions ni des
recommandations.

\paragraph{Audit.} L'audit a trouvé un défaut de code dans la v0.7 : une variable locale de la fermeture écrasait
l'amortissement du point fixe du modèle. Le choix d'un terme d'amortissement de 0,3 en v0.7 en dépendait ; il est
invalidé. Deux termes non estimés de la fermeture produisaient des oscillations de période deux ans ; ils sont
supprimés. Six limites de la v0.7 sont analysées (mécanisme, hypothèses implicites, nature, corrections possibles).

\paragraph{Transmission vers le potentiel.} La fermeture est réécrite avec deux paramètres : ρ, le rappel de la
production vers le potentiel, et θ, la part de la hausse du potentiel anticipée par la demande. Ils sont contraints par
la durée observée des écarts de production (2,3 à 2,8 ans, intervalle 1,6-5,3, AMECO et FMI 1982-2019) et par la
réaction de la consommation aux chocs permanents de revenu. Aucune estimation ponctuelle robuste n'existe : on publie
une plage (θ de 0,5 à 1,0, ρ de 0,1 à 0,3) et une valeur centrale fixée avant de regarder les stratégies (θ = 0,75 ;
ρ = 0,2). Aucune asymétrie ni dépendance à l'état n'est détectable.

\paragraph{Backtest.} Les équations sont réestimées à chaque date d'origine avec les seules données alors disponibles.
Sur 2010-2019, l'erreur sur le PIB est de 1 point (moyenne quadratique), et le modèle ne bat la prévision naïve que sur
les origines 2012 et 2016. L'emploi est sous-estimé de 1,9 \% et la dette surestimée de 20 points en 2019 : la
productivité estimée en temps réel était trop optimiste, et la charge d'intérêts suit trop lentement la baisse des
taux. Sur 2021-2025, l'inflation est sous-estimée d'environ 6 points en niveau. Les effets simulés d'une baisse de
cotisations (132 000 à 159 000 emplois pour 20\md), d'une hausse de TVA et d'une baisse d'IS sont dans les fourchettes
des évaluations ; ceux de la prime d'activité et du RSA sur l'offre de travail sont dans le haut.

\paragraph{Distribution dynamique.} Le nouveau module M16 suit le revenu réel de 41 000 ménages par huit canaux
(prélèvements et prestations, salaires, emploi, patrimoine, successions, indexation, prix, plus un contrôle de cohérence
avec le compte macroéconomique), à cinq horizons, par décile, âge, statut, catégorie sociale, configuration familiale et
statut d'occupation du logement. Trois défauts de représentation des gagnants et des perdants y ont été trouvés et
corrigés pendant la phase (âge légal, allocations chômage, cohérence nominale).

\paragraph{Stratégies S1 à S5.} Les horizons 2027-2035 sont relativement robustes ; au-delà de 2040, l'incertitude domine.
\begin{itemize}
\item \textbf{S1 (protection)} : chômage $-1{,}2$ point et pauvreté $-4{,}3$ points en 2035, aucun perdant ; dette
neutre en 2035 mais $+10$ points en 2047 (de $-13$ à $+25$ selon les réactions).
\item \textbf{S2 (emploi)} : PIB $+1{,}3$ \% en 2035, $+1{,}8$ \% en 2047 ; dette $+26$ points en 2047 (19 à 26) ; les
chômeurs perdent (allocations).
\item \textbf{S3 (productivité)} : le plus fort effet sur le PIB ($+4{,}3$ \% en 2047, 3,5 à 5,3), mais la dette augmente
de 53 points (44 à 72), le chômage de 0,9 point en 2047.
\item \textbf{S4 (redressement)} : seule stratégie qui réduit nettement la dette ($-15$ points en 2035, $-59$ en 2047,
de $-44$ à $-72$ selon la prime de risque), au prix de 0,5 \% de PIB et d'une perte de revenu réel pour tous les
déciles (1,6 à 4,8 \%) ; la pauvreté augmente de 1,4 point.
\item \textbf{S5 (global)} : effets modérés ; dette $+5$ points en 2047 (3 à 7), PIB $+0{,}9$ \% ; les déciles 3 à 6
gagnent, le premier et le dernier perdent.
\end{itemize}
Aucune stratégie ne domine les autres sur tous les critères : le rapport ne désigne pas de gagnante.

\paragraph{Ce qui change par rapport à la v0.7.} La fermeture estimée retire aux réformes d'offre l'essentiel de leurs
gains budgétaires : le relèvement de l'âge légal à 65 ans ne réduit plus la dette de 2047 que de 2 points (8,5 en v0.7),
et S3 l'augmente de 53 points (18 en v0.7). Quatre conclusions de la v0.7 sont invalidées, dont celle d'une baisse
robuste de la dette par S5. Les compositions optimisées en v0.7 pour S1, S2 et S3 ne respectent plus leur contrainte de
dette.

\paragraph{Incertitude.} Cinq sources sont distinguées : \emph{statistique} (erreurs d'estimation, durée des écarts de
production), \emph{paramétrique} (comportements calibrés : configurations faible et forte, sensibilité un à un sur
douze paramètres), \emph{structurelle} (forme de la fermeture, incidence, prime de risque souveraine — cette dernière pèse
autant que tous les comportements réunis sur la dette de 2047), \emph{chocs futurs} (environnements de la phase 3) et
\emph{mise en œuvre}. Les limites restantes sont classées au chapitre~\ref{ch:limites} ; une revue critique conclut le
rapport.
